\section{Агентная нагрузка}
\label{sec:workload}

Расчёты опираются на журналы 13\,175 агентских сессий за год, с ноября 2025 по октябрь 2026 года, в Claude Code,
Codex и Muse: 1\,528\,194 вызова модели, в основном разработка и проектирование кодовых баз. По всем вызовам средняя
глубина контекста 164 тысячи токенов (медиана 135 тысяч), на вызов в среднем приходится 1\,065 новых токенов и 512
токенов ответа, в которые рассуждения модели уже входят. Для расчётов движок генерирует по этим распределениям 40
тысяч вызовов; средние и вес крупных порций в выборке совпадают с журналами.

\textbf{Вызов, сессия и доля открытой модели.} Вызов означает один запрос агента к модели: чтение новой порции и
ответ. Сессия объединяет вызовы одного разговора агента, её контекст только растёт. Стоимость мы считаем на 1000
вызовов, потому что большая часть расходов приходится на повторное чтение глубокого контекста, а в цене за токен этого
не видно. Долей открытой модели $s$ мы называем часть вызовов пользователя, которую обслуживает открытая модель;
остальные идут к более сильной облачной модели. Модель выбирают для сессии целиком, иначе контекст пришлось бы читать
заново. Основной диапазон расчётов: $s$ от 20 до 50\%. В наших журналах на моделях не сильнее Qwen3.8-27B по индексу
Artificial Analysis \cite{aaindex} прошло не больше 17\% вызовов, но только 5\% счёта: основные деньги приходятся на
передовые модели.

\textbf{Prefill и decode.} Prefill упирается в вычисления: каждый новый токен проходит через все слои и смотрит на
весь предыдущий контекст. Decode упирается в полосу памяти: на каждом токене он читает активные веса и весь кэш. У
машин с унифицированной памятью полоса велика относительно вычислителей, поэтому decode у них идёт приемлемо, а
prefill на глубине медленно. Таблица~\ref{tab:m3max} показывает это на M3 Max с 40-ядерным GPU и 128~ГБ памяти,
модель Qwen3.8-27B в квантизации Q4\_K\_M, кэш внимания в 8 битах.

\begin{table}[t]
\centering
\caption{Замер на M3 Max (128~ГБ), Qwen3.8-27B Q4\_K\_M, llama.cpp на Metal. Скорость prefill измерена на отрезках
контекста, глубина указана по середине отрезка. Decode: 7.5--8.8 токена/с на глубине 16 тысяч и около 5 на глубине
249 тысяч; четыре параллельных потока на глубине 32 тысячи дают 13.2 токена/с суммарно по времени запросов
(14.05 по интервалам генерации) против 8.44 у одного.}
\label{tab:m3max}
\footnotesize
\begin{tabular}{lrrrrrr}
\toprule
Глубина контекста, тыс. токенов & 8 & 39 & 81 & 131 & 181 & 231 \\
Prefill, токенов/с & 173.6 & 93.7 & 49.9 & 41.2 & 27.7 & 22.2 \\
\bottomrule
\end{tabular}
\end{table}

По скоростям этих отрезков порция в 10 тысяч токенов на глубине 131--181 тысячи читается от 4 до 6 минут, а ответ
в 512 токенов генерируется за 75--93 секунды: на глубине 150 тысяч замер дал от 5.5 до 6.8 токена в секунду. На всём
профиле вызовов чтение на M3 Max занимает в среднем 28 секунд на вызов; у половины вызовов меньше 4 секунд, у десятой
доли больше 77 секунд. Ожидание сосредоточено в крупных порциях, и именно их схема переносит на сервер.
